\subsection{modèles simplifiés QG et SQG}
\textbf{Introduire QG : dynamique lente, simplifications, avantages et inconvénients sur le plan numérique. Instabilités (barotropes et baroclines) et turbulence. Puis Modèle SQG (le seul à décrire la température en surface d'une couche}
\subsection{Couplage océan-atmosphère}
\textbf{Se baser sur \parencite{Vic2024}. Influence des couplages mécaniques et thermiques, dans le cadre d'un modèle d'Eady. Incidemment, on utilisera ces couplages. Introduit d'abord un peu le couplage océan-atmosphère : Eddy-killing, pompage d'Ekman, effet d'une atmosphère sur les écoulements méso-subméso-échelle}
L'océan et l'atmosphère échangent à la fois de la chaleur et de la quantité de mouvement. \parencite{Vic2024} a étudié le couplage océan-atmosphère pour le problème d'Eady \textbf{référence vers travail original} dans le cadre d'un modèle SQG 2 couches. Le couplage thermique est représenté sous forme d'un flux turbulent de chaleur latente : 
\begin{align}
	\frac{D_h \theta_a}{Dt} = C_o(\theta_o - \theta_o) + F_a \\
	\frac{D_h \theta_o}{Dt} = C_a(\theta_a - \theta_o) + F_b
\end{align}

Couplage mécanique : 
\begin{align}
	\frac{D_h \theta_a}{Dt}+w_a N_o =F_a \\
	\frac{D_h \theta_o}{Dt}+w_o N_o =F_o \, \mathrm{et} \\
	w_o = \frac{1}{\rho_o f_0} \Vec{rot}\Vec{\tau_o}_z = \frac{1}{\rho_o f_0} (\partial_x \tau_2^y - \partial\tau_2^x)\\
	\Vec{\tau_a} = - \vec{\tau_o} \Rightarrow w_a = -\frac{\rho_o}{\rho_a}w_a
\end{align}
Où $N_o$ et $N_a$ sont les fréquences de Brunt-Vaisala \textbf{A voir comment ça se passe pour le modèle TQG, car on fait des moyennes verticales !}. En général, la friction du vent est paramétrisée de la manière suivante : 
\begin{equation}
	\Vec{\tau_o} = \rho_a C_d \abs{ \Vec{u_a} - \Vec{u_o} }(\Vec{u_o} - \Vec{u_a})
\end{equation}
Si la vitesse du vent est bien supérieure à celle du courant de surface ($\Vec{u_a} \gg \Vec{u_o}$), la relation précédente peut être linéarisée, de telle sorte que : 
\begin{equation}
	\Vec{\tau_o} = D_a (\Vec{u_o} - \Vec{u_a})
\end{equation}
\textbf{Effet des couplages thermiques/mécaniques}

\subsection{Modèle Thermal Shallow Water}
\textbf{Expliquer les équations TRSW (\parencite{Warneford2013}), justifier les termes de pression. Equilibre Geostrophique (cf \parencite{Gouzien2017})). Leur intéret, les invariants (pseudo-énergie, PV). Eventuellement les re-dériver en annexe. Après, la formulation varie selon les auteurs. Stabilité des écoulements}
\begin{align}
	\dT{h}+\div _\mathrm{h}{h\Vec{u}}=0 \\
	\dT{\Theta}+\Vec{u}\cdot\grad\Theta = - \kappa(h\Theta - H_0 \Theta_0) \\
	\dT{\Vec{u}}+(\Vec{u}\cdot \grad)\Vec{u} +f\Vec{z}\VecProd\ \Vec{u}=-\grad h\Theta+\frac{1}{2}h\grad \Theta
\end{align}
\subsection{Modèle Thermal Quasi-Geostrophic (TQG)}
\textbf{Re-dériver le modèle TQG par expansion en puissance du nombre de Rossby (si c'est trop long, mettre en Annexe). Caractéristiques du modèle (Invariants, solutions). Exemple d'utilisation (cf \parencite{Carton2025}). La formulation étant différente selon les auteurs, on se basera sur celle de \cite{Warneford2013}. }
Formulation adoptée (cas d'un terrain plat, j'ai pas trouvé de dérivation claire si la bathymétrie est variable) : 
\begin{align}
	\dT{q}+\mathrm{J}(\psi, q) = \frac{1}{R_d^2}\mathrm{J}(\psi, \theta) \\
	\dT{\theta}+\mathrm{J}(\psi, \theta) = -\lambda(\theta+\psi) \\
	q = \Delta_{\mathrm{h}} - \frac{\psi - \theta}{R_d^2} + f_0 + \beta y
\end{align}
Où $\lambda$ est un terme de refroidissement, qu'on prendra nul.

\subsubsection{Preuve (adaptée de mes notes perso et de \cite{Warneford2013})}
 On utilise le scaling suivant : 
 \begin{align}
 	x = L\tilde{x} \\
 	u = U\tilde{u} \\
 	t = \frac{L}{U} \tilde{t}
 \end{align} 
Pour la flottabilité $\Theta$ et la hauteur de couche $h$, l'équilibre géostrophique donne $f_0 U \approx \frac{\eta\Theta_0}{L}$, de telle sorte que 
\begin{align}
	h\approx H_0(1+\frac{f_0 U L}{\Theta_0 H_0}\tilde{eta}) \\
	\Theta \approx \Theta_0 (1+2\frac{f_0 U L}{\Theta_0 H_0}\theta)
\end{align}
Les équations TRSW s'écrivent donc ( pour $\lambda =0$ ) : 
\begin{align}
	\mathrm{Ro}\[ \frac{\partial\tilde{\eta}}{\partial \tilde{t}} + \tilde{\div}\tilde{\eta}\tilde{\Vec{u}} \] +\mathrm{Bu}\cdot\tilde{\div} \tilde{\Vec{u}} =0 \\
	\frac{\partial \tilde{\theta}}{\partial \tilde{t}} + \tilde{\Vec{u}}\cdot \grad \tilde{\theta} = 0 \\
	\mathrm{Ro}\[ \frac{\partial \tilde{\Vec{u}}}{\partial \tilde{t}} +(\tilde{\Vec{u}}\cdot \tilde{\grad})\tilde{\Vec{u}}\]+ (1+\tilde{\beta}\tilde{y})\Vec{z}\VecProd\tilde{\Vec{u}} = - \tilde{\grad}{\tilde{\eta}+\tilde{\theta}}-\frac{\mathrm{Ro}}{\mathrm{Bu}}\[ 2\tilde{\theta}\tilde{\grad}\tilde{\eta}+\tilde{\eta}\tilde{\grad}\tilde{\theta}\]
\end{align}
Avec $\mathrm{Ro} = \frac{U}{f_0L}$, $\mathrm{Bu}=\frac{\Theta_0 H_0}{f_0 L^2} =(\frac{R_d}{L})^2$, $\tilde{\beta} = \frac{\beta L}{f_0}$. Par la suite, on arrêtera d'utiliser les $\tilde{}$, pour ne pas alourdir les notations. Sachant que $\mathrm{Ro} \ll 1$, on peut développer le champ de vitesse en puissances de $\mathrm{Ro}$ : 
\begin{equation}
	\Vec{u}= \Vec{u_0} + \mathrm{Ro}\Vec{u_1} +\mathrm{Ro}^2\Vec{u_2} + ...
\end{equation}
On suppose aussi que $\tilde{\beta} \approx \mathrm{Ro}$. A l'ordre zéro, on a 
\begin{align}
	\div \Vec{u_0} = 0 \\
	\dT{\theta} + \Vec{u_0}\cdot \grad\theta = 0 \\
	\Vec{z}\VecProd \Vec{u_0} = - \grad{\eta+\theta}
\end{align}
La première relation implique l'existence d'une fonction de courant $\psi$ telle que $u_0 = -\frac{\partial \psi}{\partial y}$ et $v_0 = \frac{\partial \psi}{\partial x}$. La troisième est simplement l'équilibre géostrophique. Afin d'obtenir une équation d'évolution pour $\Vec{u_0}$, on passe à l'ordre 1, ce qui donne 
\begin{align}
	\dT{\eta}+\div{\eta\Vec{u_0}} +\mathrm{Bu}\cdot \div{\Vec{u_1}} = 0 \\
	\Vec{u_1}\cdot\grad \theta  = \frac{\theta\eta}{\mathrm{Bu}} \\
	\dT{\Vec{u_0}}+(\Vec{u_0}\grad)\Vec{u_0} + \Vec{z}\VecProd \Vec{u_1} +\beta y \Vec{z}\VecProd\Vec{u_0} = \frac{1}{\mathrm{Bu}}\[ 2\theta \grad \theta +\eta\grad\theta\]\label{eq:dT_u0}
\end{align}
Pour obtenir la vorticité absolue, il suffit de prendre le rotationnel de l'équation \ref{eq:dT_u0}. En utilisant les relations $\Vec{mathrm{rot}} f \Vec{A} = \grad f\VecProd \Vec{A} + f \Vec{\mathrm{rot}}\Vec{A}$ et $\Vec{\mathrm{rot}} = \Vec{z}\div \Vec{A} - \dZ{Vec{A}}$. On a donc, en projetant selon $\Vec{z}$ : 
\begin{equation}
	\dT{\eta_a} +\div Vec{u_1 } = -\frac{1}{\mathrm{Bu}}\(\frac{\partial \theta}{\partial x}\frac{\partial \eta}{\partial y}-\frac{\partial \theta}{\partial xy}\frac{\partial \eta}{\partial x}\)
\end{equation}
$\Vec{u_1}$ s'élimine simplement puisque $\div\Vec{u_1} = -\frac{1}{\mathrm{Bu}} ( \dT{\eta} +\Vec{u_0}\cdot\grad\eta )$, et d'après l'équilibre géostrophique, $\eta = \psi - \theta$. On en déduit donc les équations TQG : 
\begin{align}
\dT{q}+\mathrm{J}(\psi, q) = \frac{1}{R_d^2}\mathrm{J}(\psi, \theta) \\
\dT{\theta}+\mathrm{J}(\psi, \theta) =0 \\
q = \Delta_{\mathrm{h}} - \frac{\psi - \theta}{R_d^2} + 1 + \tilde{\beta} y\\
\end{align}
Où $\mathrm{J}(A, B) = \partial_x A \partial_y B - \partial_y A \partial_x B$ est le Jacobien 2D et $\Delta_{\mathrm{h}}= \partial_x^2+\partial_y^2$ le Laplacien 2D
Propriétés (démontre celles qui sont simples) : 
\begin{itemize}
	\item Conservation de la pseudo-énergie (du moins sans dissipation)
	\item Conservation de la PV intégrée (\textbf{si le contour délimitant la surface est une isotherme})
\end{itemize}
